\origpage{40}{ROBERTSON AND SEYMOUR}
{\centering
{\Large Graph Minors. IX. Disjoint Crossed Paths\par}\medskip

Neil Robertson$^*$\par
\textit{Department of Mathematics, Ohio State University, 231 West 18th Avenue, Columbus, Ohio 43210}\par
and\par
P. D. Seymour\par
\textit{Bellcore, 445 South St., Morristown, New Jersey 07960}\par
Received February 10, 1987; revised July 12, 1989\par
Let $G$ be a graph, with a cyclic order imposed on a subset of its vertices (called ``special'' vertices). We show that either\par
\hspace*{12pt}(i) modulo ($\le 3$)-separations, $G$ can be drawn in a disc with no crossings except in one ``small'' area, and with its special vertices on the outside in the correct order, or

\hspace*{12pt}(ii) there is a partition of the special vertices into two ``semicircles,'' and there is a large collection of vertex-disjoint paths of $G$ running from one semicircle to the other, such that each of these paths is either crossed by another or lies between two others.

This is basically a lemma to be used in later papers.\par
}

\sectionhead{1. Introduction}

We hope to prove in a later paper that for any graph $H$, every graph with no minor isomorphic to $H$ can be constructed by piecing together in a tree structure graphs which can ``almost'' be drawn in surfaces of bounded genus. This result appears to have a number of important consequences, including the truth of Wagner's well-quasi-ordering conjecture for graph minors, and the existence of a polynomial algorithm for the Disjoint Connecting Paths problem. However, the proof of this embedding structure theorem is long, and it is convenient to present it in sections. Here we%
\origfootnote{This research was performed under a consulting agreement with Bellcore.}
